\section{Source-layer model-set input}\label{sec:superspace}

This appendix records the only cut-and-project facts used in the main text. They justify the benchmark source law and the existence of well-defined patch frequencies; all theorem-level arguments in the body of the paper are symbolic and arithmetic.

Let $E=\RR^d$, $H=\RR^n$, and let $\Gamma\subset E\times H$ be a lattice such that $\pi_E|_{\Gamma}$ is injective and $\pi_H(\Gamma)$ is dense. Writing $\Gamma_E:=\pi_E(\Gamma)$ and $x^\star:=\pi_H(\gamma)$ for the unique $\gamma\in\Gamma$ with $\pi_E(\gamma)=x$, a compact window $W\subset H$ with nonempty interior and Haar-null boundary defines the regular model set
\[
\Lambda(W,u):=\{x\in\Gamma_E:\ x^\star\in W+u\}\subset E
\]
at phase $u\in H$. Varying $u$ moves along the hull, and the hull has $(E\times H)/\Gamma$ as its natural equicontinuous factor. In particular, a one-dimensional scan is obtained by coding a torus orbit through a finite observation map \cite{Moody2000,BaakeGrimm2013}.

\begin{proposition}[Source-layer input from regular model sets]\label{prop:unique-ergodicity}
Let $\Lambda(W,u)$ be a regular model set as above, and let $X_{\Lambda}(u)$ be its hull in the local topology with the $E$-translation action. Then $(X_{\Lambda}(u),E)$ is uniquely ergodic. In particular, for $\mu_H$-a.e.\ phase $u$ and every finite patch $P$, the frequency $\mathrm{freq}(P;\Lambda(W,u))$ exists and agrees with the unique invariant patch frequency. The scan dynamics can therefore be modeled by a torus translation with a well-defined symbolic factor law \cite{Moody2000,BaakeGrimm2013}.
\end{proposition}

This is the only place where source-layer cut-and-project input is required. Once that input is fixed, the main argument reduces to the symbolic-arithmetic stabilization theory developed in the body of the paper.
